\documentclass[11pt,a4paper]{article}
\usepackage[utf8]{inputenc}
\usepackage[T1]{fontenc}
\usepackage[polish]{babel}
\usepackage{lmodern}
\usepackage{geometry}
\usepackage[hidelinks]{hyperref}
\geometry{margin=27mm}
\title{Objawy lękowe a elastyczność poznawcza i podejmowanie decyzji}
\author{}
\date{Prospekt badawczy, wersja 0.1.0, 22 listopada 2026}
\begin{document}
\maketitle
\begin{abstract}
Projekt MSC-ANX-P01 ma status \texttt{IDEA\_REVIEW}. Ocenić, czy nasilenie objawów lękowych wiąże się z kosztami przełączania reguł i zmianą strategii decyzyjnych przy niepewności. Dokument zawiera planowane hipotezy i metody, a nie wyniki empiryczne.
\end{abstract}
\section{Problem i cel badania}
Ocenić, czy nasilenie objawów lękowych wiąże się z kosztami przełączania reguł i zmianą strategii decyzyjnych przy niepewności.
\section{Hipotezy robocze}
Wyższe nasilenie objawów lękowych będzie związane z większym kosztem przełączania oraz silniejszą awersją do niepewności; hipotezy nie przesądzają o kierunku zależności przyczynowej.
\section{Projekt badania}
Przekrojowe badanie online osób pełnoletnich. Pomiar objawów lękowych za pomocą zwalidowanej wersji GAD-7; komputerowe zadanie przełączania reguł i zadanie wyboru probabilistycznego, z próbami kontrolnymi. Warunki uczestnictwa, kryteria wykluczenia, uzasadnienie liczebności próby i procedura świadomej zgody zostaną ustalone przed prerejestracją.
\section{Operacjonalizacja}
Pierwszorzędowy: koszt przełączania (różnica mediany czasu poprawnych reakcji w próbach switch versus repeat). Drugorzędowe: trafność wyboru, odsetek decyzji ryzykownych i modelowany parametr wrażliwości na niepewność.
\section{Plan analiz}
Regresja kosztu przełączania na wynik GAD-7 z korektą wieku, snu i objawów depresyjnych; osobno hierarchiczny model wyborów probabilistycznych. Korekta wielokrotnych porównań dla miar drugorzędowych. Przed analizą danych należy zamrozić estymand, reguły czyszczenia danych, plan braków danych i kryteria analiz wrażliwości.
\section{Etyka, bezpieczeństwo i dane}
Nie rozpoczęto rekrutacji. Konieczne jest rozstrzygnięcie wymogów opinii komisji etycznej, minimalizacji danych, retencji i informowania o dobrowolności udziału. Nie zbierać danych pozwalających na identyfikację w publicznym repozytorium.
\section{Wyniki}
Brak danych i wyników. Sekcja zostanie uzupełniona dopiero po zrealizowaniu zatwierdzonego protokołu.
\section{Ograniczenia}
Samoopis nie jest diagnozą; nakładanie się lęku i depresji, efekt szybkość--trafność, wpływ urządzeń i rekrutacja ochotnicza ograniczają uogólnianie.
\section{Dyskusja}
Interpretacja będzie zależała od trafności narzędzi, jakości danych i alternatywnych wyjaśnień. Projekt nie zakłada z góry potwierdzenia hipotez.
\section*{Literatura}
Weryfikacja prac źródłowych i przygotowanie bibliografii z kontrolą DOI należą do kolejnego etapu. Nie umieszczono niezweryfikowanych cytowań.
\end{document}
